\documentclass[../../main.tex]{subfiles}
\begin{document}

% 年度の大きな表示
\begin{center}
  {\fontsize{48pt}{58pt}\selectfont \textbf{1999}\normalsize\textbf{年度}}
\end{center}

\vspace{10mm}

% 本文


\vspace{15mm}

% 出題テーマと難易度の表
\noindent\textbf{出題テーマと難易度}

\vspace{5mm}

\begin{center}
  \shadowbox{%
    \begin{tabular}{|c|p{8cm}|c|}
      \hline
      \textbf{問題番号} & \textbf{テーマ} & \textbf{難易度} \\
      \hline
      第1問 & 不等式の大小比較（凸関数と冪乗） & \\
      \hline
      第2問 & 空間図形（正$n$角錐の体積の最大化） & \\
      \hline
      第3問 & 図形（三角形を2等分する線分の最小値） & \\
      \hline
      第4問 & 定積分と階乗の等式の証明 & \\
      \hline
      第5問 & 複素数列の漸化式と極限点 & \\
      \hline
    \end{tabular}%
  }
\end{center}

\vspace{15mm}

% 解答
\noindent\textbf{解答}

\vspace{5mm}

\begin{center}
  \shadowbox{%
    \renewcommand{\arraystretch}{1.6}
    \begin{tabular}{|c|c|p{8cm}|}
      \hline
      \textbf{問題番号} & \textbf{小問} & \textbf{答え} \\
      \hline
      第1問 & - & $p=1$または$a=b$のとき$A=B$，$0<p<1$かつ$a\ne b$のとき$A>B$，$1<p$かつ$a\ne b$のとき$A<B$ \\
      \hline
      \multirow{2}{*}{第2問} & (1) & $V_n=\dfrac{2\sqrt{3}}{27}\,n\sin\dfrac{\pi}{n}\cos\dfrac{\pi}{n}$ \\
      \cline{2-3}
                             & (2) & $\dfrac{2\sqrt{3}\,\pi}{27}$ \\
      \hline
      第3問 & - & $\min PQ=\dfrac{\sqrt{2}}{2}a\ (0<a\le1)$，$\sqrt{-\dfrac{1}{2}a^2+a}\ (1\le a<2)$ \\
      \hline
      第4問 & - & 証明問題 \\
      \hline
      \multirow{2}{*}{第5問} & (1) & $P_\infty=\dfrac{3}{7}\left\{(2-\sqrt{3})+(2+\sqrt{3})i\right\}$ \\
      \cline{2-3}
                             & (2) & $z=\dfrac{3}{28}\left\{(13-5\sqrt{3})+(3+3\sqrt{3})i\right\}$ \\
      \hline
    \end{tabular}%
    \renewcommand{\arraystretch}{1.0}
  }
\end{center}

\vspace{15mm}

% 解答時間
\noindent\textbf{試験形態}

\vspace{5mm}

\begin{center}
  \shadowbox{%
    \begin{tabular}{p{8cm}}
      （要確認）
    \end{tabular}%
  }
\end{center}

\end{document}
